\section{Experimental Environment and Data}

\subsection{Official data source}

The experiment uses \texttt{aloeL.jpg} and \texttt{aloeR.jpg} from the \texttt{4.x/samples/data} directory of the official OpenCV repository, \texttt{opencv/opencv}. This image pair was selected because it is an official OpenCV sample with an explicit left--right relationship, and the official Python stereo-matching example also provides a reproducible SGBM processing strategy for the aloe images. It therefore supports the course experiment without introducing an additional dataset.

Both original JPG images have dimensions of $1282\times1110$ and contain three channels. The program synchronously applies \texttt{pyrDown} once to the left and right images, giving an actual matching size of $641\times555$. \texttt{aloeGT.png} in the same official directory is an optional disparity reference image; this experiment does not treat it as a calibration parameter, a depth scale, or an error ground truth.

\begin{table}[H]
  \centering
  \caption{Experimental data and software environment}
  \label{tab:environment}
  \begin{tabular}{ll}
    \toprule
    Item & Measured version or description \\
    \midrule
    Data directory & Official OpenCV \texttt{4.x/samples/data} \\
    Input files & \texttt{aloeL.jpg}, \texttt{aloeR.jpg} \\
    Original dimensions & $1282\times1110$, three channels \\
    Python & 3.13.5 \\
    OpenCV runtime & 5.0.0 \\
    NumPy & 2.5.3 \\
    \bottomrule
  \end{tabular}
\end{table}

Version information, input paths, and run parameters come from the program-generated \codepath{task2/results/params.json}. The SHA-256 values and download URLs of the input files are recorded in \codepath{task2/results/run_evidence.md}, allowing the experiment to be checked in the same environment without relying on image statistics copied manually into the report.

\subsection{Actual run command}

The results in this report were generated by the following default command. The run did not use \texttt{--no-pyrdown} and did not apply any additional parameter adjustment:

\begin{lstlisting}[language=bash]
task2/.venv/bin/python task2/src/stereo_depth.py \\
  --left task2/data/aloeL.jpg \\
  --right task2/data/aloeR.jpg \\
  --out-dir task2/results
\end{lstlisting}

The program first checks that both images can be read and that their dimensions agree. It then performs grayscale conversion, synchronized downsampling, StereoSGBM matching, disparity filtering, and relative-depth computation. No calibration file from another image pair is used; all results come from the outputs in the root directory corresponding to the command above.
